\section{Introduction}
\label{sec:intro}

Three-dimensional fluorescence imaging normally records a stack of focal planes, moving either the specimen or the focus. Fixed-focus alternatives encode depth into planar images and recover the volume computationally. Most encode depth on the detection side, giving each depth a distinct point-spread function with a lenslet array, diffuser or phase mask \cite{broxton2013,antipa2018,yanny2020,guo2019,xue2020}; others encode it in the excitation with speckle or patterned illumination, usually by translating or regenerating the pattern \cite{guo2016,pascucci2019,yeh2019}. Each choice fixes what has to move or be replaced between measurements and how much depth information a planar image carries.

We study an excitation-side encoder in which nothing moves except the direction of the incident beam. A plane wave is scanned in angle across one fixed phase-only plate placed a short distance before an intermediate illumination image plane; a high-NA relay maps the resulting field into the specimen; and a conventional low-NA camera at a fixed focus records one image per angle (\cref{fig:system}). The plate converts each input direction into a three-dimensional excitation pattern, and because a tilt translates the diffracted pattern by an amount that grows with propagation distance, different depths see different sequences of excitation across the scan. This gives depth-distinguishing measurements without axial scanning, refocusing, a moving diffuser or a reconfigurable modulator.

Turning these measurements into volumes needs a forward model that is accurate at the scale of the illumination structure. We simulate the excitation with a vectorial high-NA model on a fine grid and integrate it exactly over each reconstruction voxel. An illumination--emission factorisation that follows from linear incoherent fluorescence makes this tractable, and the full operator builds in about two minutes on one GPU. We reconstruct with two standard solvers, box-constrained total-variation least squares and multi-slice plug-and-play ADMM with one pretrained 2D denoiser on three orthogonal slice sets \cite{venkatakrishnan2013,majee2021,zhang2017}; the encoder comparison holds with both.

A second question arises once volumes look convincing. A reconstruction may show two locations joined by fluorescent structure even when the measurements also admit a volume in which they are disjoint, so the prior, not the data, made the choice. We therefore complement the imaging study with a measurement-level analysis of connectivity: for a query, we either certify that every volume consistent with the images and a calibrated noise model agrees, or exhibit two consistent volumes with opposite connectivity.

\paragraph{Contributions.}
\begin{enumerate}\itemsep0pt
\item \textbf{Acquisition.} A fixed-focus 3D fluorescence acquisition in which only the illumination angle changes through one fixed, defocused phase encoder. We explain its depth encoding by the depth-dependent shear of the excitation with input tilt, measure that shear numerically, and compare it under matched photon budgets with a no-plate angle scan, an in-focus amplitude-pattern (DMD) encoder and a defocused amplitude mask, all through the same vectorial relay, detector and solvers, over several encoder realisations and plate defocus distances.
\item \textbf{Forward model and reconstruction.} An exactly voxel-integrated operator built from a fine-grid vectorial model through an illumination--emission factorisation, which reduces the forward discrepancy of direct coarse-grid simulation from 25.7\% to 1.07\%, combined with standard regularised recovery. On mouse neuron volumes from six cells of three animals, dominant dendrites and their depth order are retained at eight-fold compression ($4{,}608$ measurements for $36{,}864$ voxels).
\item \textbf{Structural reliability.} A query-level analysis of which reconstructed connections the measurements establish. On a validated reduced field at the system's operating noise level, many connectivity queries are not determined by the data, even when the reconstruction is correct. Using every available direction at a fixed photon budget resolves none of them, a calibrated regularity prior does not remove the ambiguity, and fine-grid model mismatch widens it.
\end{enumerate}

